% ============================================================================
%  Resolución consolidada de observaciones (Artículo 46° y 47.1°, FEP01.26)
%  para la Oferta Final (\Instancia{final}).
%  Consolida el 100 % de las observaciones de la Revisión Final del Informe 2
%  (60 observaciones de Revisión 1 + 6 observaciones complementarias de
%  Revisión 2 = 66 observaciones). El Informe 3 económico no recibió informe
%  escrito de observaciones; su retroalimentación oral fue incorporada en la
%  versión económica 3.2 (Bitácora A-6 v3.2).
% ============================================================================
\begin{tablaObservaciones}{Informe 2 (Revisión Final) e Informe 3}

\grupoObservaciones{Presentación y control transversal}
\observacion[sub=0]{GEN-01}{Las tablas de respuesta ocupan una sección que altera el índice obligatorio de los subdocumentos. (rev. final, p. 3)}{La tabla de respuestas se consolida en este anexo independiente del Art. 46°. Los subdocumentos S1, S2 y S9 conservan intacta la estructura de títulos literales de la Aclaración N.° 2.}{Anexo Art. 46° y aperturas de S1, S2 y S9. Estado: Incorporada}
\observacion[sub=0]{GEN-02}{Los folios reinician por archivo. La revisión exige numeración correlativa del conjunto para la oferta final. (rev. final, pp. 1--3, 14, 48)}{Configuración de instancia actualizada a \texttt{final} con \texttt{ModoFolio\{continuo\}} para el ensamblaje definitivo del Sobre N.° 2 y conciliación con la hoja resumen de cada sobre.}{configuracion/instancia.tex y hojas resumen de sobres finales. Estado: En cierre de ensamblaje final}
\observacion[sub=0]{GEN-03}{S1, S2 y T-6 reiteran cláusulas defensivas. S1 mantiene prosa institucional de folleto. (rev. final, pp. 3--5, 7--8, 48)}{Se depuró la presentación de capacidades, compromisos y supuestos. S1 retiró misión, visión y valores y eliminó promesas genéricas no respaldadas en v3.2. T-6 remite al análisis de capacidad de S1 y al Sobre N.° 1.}{S1 §§1.1 y 1.4, S2 §§2.1--2.5 y T-6. Estado: Incorporada}
\observacion[sub=0]{GEN-04}{Las declaraciones de S3, S4 y S6 tienen seis columnas y la revisión las observa por el límite general del cuerpo. (rev. final, pp. 11, 14, 23, 48)}{Se conservan los seis campos expresamente exigidos por el Comunicado 10 §7.2 distinguiendo la declaración obligatoria de las tablas analíticas del cuerpo, con criterio uniforme en S3, S4 y S6.}{Declaraciones de uso de IA de S3, S4 y S6. Estado: Criterio documentado (pendiente validación en S3/S4/S6)}
\observacion[sub=0]{GEN-05}{S6, S8 y S9 declaran uso alto de IA con revisión concentrada. Se solicita especial atención a la simulación de riesgos y al diseño de pruebas. (rev. final, pp. 23, 31--32)}{Se documentó la revisión y trazabilidad de S9/T-13/T-17, se actualizaron S1, S2 y S9, y la declaración de S8 nombra tres revisores por sección indicando qué verificó cada uno, incluida la simulación sobre la malla conciliada.}{Declaraciones de IA de S1, S2, S8 y S9 (incorporado); S6 verificado por D3. Estado: Incorporada}
\observacion[sub=0]{GEN-06}{El T-11 carece de declaración propia y los T-14, T-15 y T-18 remiten a la del capítulo. (rev. final, pp. 17--18, 27, 48)}{T-11 cierra con Referencias y Declaración de uso de IA propias; T-16 y T-19 incluyen también declaración propia. T-14, T-15 y T-18 fueron integrados por D2 en el cierre de planificación.}{Formularios T-11, T-16, T-19 y cierre D2 (T-14, T-15, T-18). Estado: Incorporada}

\grupoObservaciones{Empresa y comprensión de la necesidad (S1, T-6 y S2)}
\observacion[sub=1]{EMP-01}{La experiencia de mayor escala declarada cubre 120 tractocamiones frente a 374 camiones y 96.000 viajes anuales del caso, sin explicar cómo se compensa la brecha. (rev. final, pp. 4--5, 39)}{S1 §1.2 y §1.4 acreditan la compensación de escala mediante el dimensionamiento real alineado con v3.2: 37,0 FTE base en M13--M17, refuerzos de puesta en producción hasta 38,50 FTE en H6/H7, pico contratado de 39,23 FTE en M7, Mesa de Ayuda N1 6+1 FTE y servicio de campo externo por visita/rondas en los cinco terminales.}{S1 §§1.2 y 1.4 y Formulario T-6. Estado: Incorporada}
\observacion[sub=1]{EMP-02}{El gerente general es también jefe de proyecto y representante legal con dedicación de 100 \% entre M1 y M20. Falta explicar la conducción de la empresa en ese período. (rev. final, p. 5)}{Se documentaron las subrogancias gerenciales y la cobertura de funciones corporativas durante la dedicación contractual de M1--M20.}{S1 §1.3 y tabla de subrogancias. Estado: Incorporada}
\observacion[sub=1]{EMP-03}{Se vinculan entidades certificadoras y auditoras reales con expedientes ficticios. (rev. final, pp. 5, 48)}{Se armonizaron las referencias e identificadores en S1 §§1.3--1.5 y Anexo 1.A (actualizado al 25 de noviembre de 2026), remitiendo el cotejo documental al Sobre N.° 1 conforme al Art. 34.1 y Art. 39.2.}{S1 §§1.3--1.5 y Anexo 1.A, sección de plan de certificación. Estado: Incorporada}
\observacion[sub=1]{EMP-04}{Las alianzas exclusivas se remiten a un capítulo y un formulario ausentes del Informe 2. (rev. final, p. 5)}{S1 §1.6 delimita explícitamente las alianzas corporativas vigentes alineadas 1:1 con v3.2 (Microsoft Azure Chile Central/Brazil South CSP, Dell, Fortinet, iWave, Teltonika, Technoton, Ground Control/Iridium, GitLab y SonarSource), radicando las cartas de respaldo en el Sobre N.° 1 y su resumen en el Capítulo 12 / Formulario T-8.}{S1 §1.6 y Formulario T-8. Estado: Incorporada}
\observacion[sub=2]{NEC-01}{La estimación de 338,15 viajes diarios de temporada alta y la cota de 450 viajes diarios dependen de un supuesto de 210 viajes diarios en valle. (rev. final, pp. 7--8)}{S2 documenta la ecuación de balance anual, la cota de estrés y su calibración obligatoria con el histórico real del sistema en H1 (M2).}{S2 §§2.1 y 2.2.3. Estado: Incorporada}
\observacion[sub=2]{NEC-02}{El diagnóstico presenta los umbrales de continuidad como compromisos propios y usa la expresión «análisis pericial». (rev. final, pp. 7--8)}{S2 atribuye expresamente a las Bases Oficiales la disponibilidad (99,9 \%), RTO ($\le 4$ h), RPO ($\le 15$ min) y autonomía (72 h), corrigiendo las citas de numerales y retirando la expresión «análisis pericial».}{S2 §§2.2.2 y 2.5.1. Estado: Incorporada}
\observacion[sub=2]{NEC-03}{La figura de cadena de valor queda aislada por espacios en blanco. (rev. final, p. 7)}{La Figura 2.1 fue redibujada en TikZ vectorial nativo a 9,5 pt integrada al flujo de S2 §2.2, recompilada y verificada sin espacios muertos ni errores en el PDF final (30 págs.).}{S2 §2.2, Figura 2.1 (PDF final verificado). Estado: Incorporada}

\grupoObservaciones{Alcance y cumplimiento técnico (S3 y T-12)}
\observacion[sub=3]{ALC-01}{SA-01 fija febrero de 2027 como inicio y el calendario depende de esa fecha para la ventana de M16. (rev. final, pp. 11--12, 46--47)}{Meses contractuales separados del calendario civil SA-01; curva física 10/60/60/52 y condiciones de firma, stock y paso por terminal explícitas.}{S3 supuestos y tabla 3-montaje; S7 calendario; T-12 SA-01 y T-18. Estado: Incorporada documentalmente}
\observacion[sub=3]{ALC-02}{La reversión se resume en una afirmación de que no se pierde información. (rev. final, p. 11)}{El procedimiento distingue preservación, reenvío, deduplicación, conciliación y retorno A/B de recuperación presencial. El ensayo CP-SYS-20 se exige antes del paso.}{S3 sec:3-reversion; T-18 tabla t18-reversion; T-17 CP-SYS-20. Estado: Incorporada en diseño; ensayo futuro}
\observacion[sub=3]{ALC-03}{El T-12 contiene 42 requerimientos propios y usa los RT como referencias, sin declarar su cumplimiento individual. (rev. final, pp. 11--12, 33, 46--47)}{Inventario individual incorporado: 374 RT transversales, 27 requisitos del Caso y 13 cruces administrativos, con localizador, estado documental y verificación prevista. La acreditación de filas sin evidencia específica se mantiene diferenciada.}{T-12 tablas t12-transversales, t12-caso y t12-administrativas. Estado: Parcial; inventario incorporado y evidencia por acreditar}
\observacion[sub=3]{ALC-04}{El T-12 declara «Sí» uniforme, incluso en RF-003, RNF-005 y RF-014, que dependen de autorizaciones o pruebas futuras. (rev. final, pp. 11--12, 46--47)}{Los 42 requisitos propios distinguen compromiso, evidencia de diseño y comprobación futura, incluidas autorizaciones, adhesión y retorno satelital. No se atribuyen resultados de ensayos aún no ejecutados.}{T-12 RF-003, RF-014 y RNF-005; S3 catálogo. Estado: Incorporada documentalmente}
\observacion[sub=3]{ALC-05}{La numeración de las tablas 3.1--3.7 del T-12 coincide con la del S3. (rev. final, pp. 11--12, 47)}{Tablas del formulario con prefijo propio T-12 y remisiones separadas de las tablas del cuerpo S3.}{T-12 numeración y referencias. Estado: Incorporada}
\observacion[sub=3]{ALC-06}{El cuerpo del S3 se refiere a versiones previas del informe. (rev. final, pp. 11--12)}{S3 se presenta como oferta autónoma; el historial se conserva en el expediente interno y esta tabla.}{S3 contenido; registro D2 de conciliación. Estado: Incorporada}

\grupoObservaciones{Arquitectura y datos (S4, T-11 y S5)}
\observacion[sub=4]{ARQ-01}{Flutter, React y Next.js no tienen versiones y soporte definidos. La versión Go declarada carece de respaldo verificable para la fecha de oferta. (rev. final, pp. 13--14, 43)}{Se alinearon Flutter 3.50.0, React 19.3.0, Next.js 16.x LTS y Go 1.27.4. Se documentaron fuentes, mantenimiento y fijación de parches.}{S4, tecnologías de software, y T-11. Estado: Incorporada}
\observacion[sub=4]{ARQ-02}{Los ocho ADR llevan una misma fecha y carecen de estado por decisión. (rev. final, pp. 14--15)}{Cada decisión registra alternativa, descartes, criterio y estado propuesto. Se distingue el registro de la oferta del acta de aprobación y la fecha efectiva.}{S4, registro de ADR. Estado: Incorporada en el diseño}
\observacion[sub=4]{ARQ-03}{La mesa de ayuda y la dotación operativa se remiten al S10 sin dimensionamiento. (rev. final, pp. 15--16, 18, 44)}{El S4 resume el dimensionamiento: 576 contactos al mes con Erlang C, dos agentes en el relevo de madrugada y uno el resto del día, 851 horas cubiertas por seis agentes más una jornada de suplencias (siete en total), y el campo con media jornada interna más un proveedor externo. El cálculo completo y su sensibilidad están en el S10.}{S4 §4.2, dimensionamiento, y S10 §10.3. Estado: Incorporada}
\observacion[sub=4]{ARQ-04}{El alcance de ExpressRoute hasta Chile Central queda sujeto a confirmación del proveedor. (rev. final, pp. 17--18, 29, 31, 46--47)}{Se nombran los proveedores de Santiago (Cirion, Equinix y PitChile). Se exige confirmación escrita antes de contratar y la conmutación a la VPN se prueba en la Etapa 1. Si el circuito no llega, se contrata Premium o dos VPN por rutas distintas. Queda como R-33 en el T-16.}{S4 §4.2, enlaces, y T-16 R-33. Estado: Incorporada}
\observacion[sub=4]{ARQ-05}{El DR de datos personales en Brazil South carece de vía habilitante documentada. R-24 depende de una actuación del mandante antes de H2 y afecta H7. (rev. final, pp. 17--20, 30--31, 45--47)}{Antes del H2, Curimón valida con asesoría legal externa el instrumento contractual con garantías adecuadas (arts. 27 letra b y 28, Ley N.º 21.719). Mientras tanto, Brazil South recibe sólo artefactos técnicos y datos sintéticos, y los datos personales quedan en Chile Central con HA entre zonas (réplica sincrónica, conmutación 60--120 s). Si no hay vía habilitante, se selecciona y ensaya un sitio secundario admisible (RPO $\le 15$ min, RTO $\le 4$ h, RT-07.04), y H7 exige ese ensayo con acta (R-24), incorporado en T-17.}{S4 §4.3, S5 privacidad, S8/T-16 R-24, T-8 y T-17 (Protocolo H-07 y CP-PERF-03). Estado: Incorporada (vía habilitante antes de H2)}
\observacion[sub=4]{ARQ-06}{La tabla del T-11 se identifica con el mismo número que la tabla de tecnologías de S4 y remite las referencias al capítulo. (rev. final, pp. 17--18, 47)}{Las tablas del T-11 se numeran T-11.n. El formulario trae sus Referencias y su ficha de fabricante por fila.}{Formulario T-11. Estado: Incorporada}
\observacion[sub=4]{ARQ-07}{Persisten brechas de cobertura real, parque de terceros, conductor sin rastro, consultas sin respuesta y autonomía de la solución satelital referida en la revisión. (rev. final, p. 18)}{Una tabla da, para cada una de las cinco brechas, el dato que falta, la prueba o consulta con su plazo y la contingencia. La autonomía ofertada (288 h del G26I y la prueba de 72 h) queda separada del producto satelital del Informe 1.}{S4 §4.2, supuestos y brechas, y T-16 R-04 y R-08. Estado: Incorporada}
\observacion[sub=4]{ARQ-08}{[Observación complementaria Rev. 2] La Figura 4.13 dibuja la «Autoridad tributaria» dentro del bloque de San Bernardo, junto al sistema contable, sin mostrar el enlace externo con el emisor ni con la capa anticorrupción. (rev. final, p. 18)}{La autoridad tributaria queda fuera del sitio. La línea al sistema contable lleva el rótulo «Internet del emisor», y el texto explica que el emisor se comunica por su propio enlace y que la capa anticorrupción solo le pide la emisión.}{S4 §4.2, Figura 4.13. Estado: Incorporada}
\observacion[sub=5]{DAT-01}{El diccionario de datos se incorporó dentro del archivo del S5. (rev. final, pp. 20, 48)}{El diccionario se entrega como PDF independiente y S5 conserva modelos, reglas y remisión al anexo.}{S5 y \texttt{AUDIT-Subdocumento5-Anexos.pdf}. Estado: Incorporada}
\observacion[sub=5]{DAT-02}{Las tablas 5.2, 5.3, 5.5 y 5.6 contienen prosa explicativa en las celdas. (rev. final, pp. 20, 48)}{Las reglas y excepciones pasan al texto y al diccionario. Las celdas comparan criterios, medidas y atributos.}{S5, secciones de almacenamiento, calidad, retención y migración. Estado: Incorporada}
\observacion[sub=5]{DAT-03}{La declaración de S5 deja una página de continuación con encabezado y sin filas. (rev. final, pp. 20, 48)}{Se recompiló la declaración con sus filas y se verificó el cierre del archivo.}{S5, declaración de uso de IA. Estado: Incorporada}
\observacion[sub=5]{DAT-04}{El 99,9 \% mensual de servicios críticos en S5 debe cotejarse con el SLA de S4. (rev. final, p. 20)}{S4, S5, S10 y T-17 (CP-PERF-03) usan los mismos cuatro servicios críticos (99,9 \%) y servicios altos (99,5 \%), con denominador mensual único y exclusión exclusiva de ventanas aprobadas (Art. 79.4).}{S4, S5, S10 §10.2 y T-17 (CP-PERF-03). Estado: Incorporada}
\observacion[sub=5]{DAT-05}{[Observación complementaria Rev. 2] La entidad TRANSPORTISTA cumple tres papeles en la Figura 5.1 con el mismo identificador de transportista sin regla cuando difieren, y el detalle de liquidación tiene clave primaria por viaje y no por asignación, con riesgo de doble liquidación ante reasignación. (rev. final, p. 20)}{Se separan titular del equipo, empleador y beneficiario. El detalle identifica la asignación ejecutada y restringe a un cargo base por tramo incluso con reasignación, con variante de prueba incorporada en T-17 (CP-SYS-10).}{S5 Figura 5.1, diccionario y T-17 (CP-SYS-10). Estado: Incorporada}

\grupoObservaciones{Metodologías (S6, T-9 y T-10)}
\observacion[sub=6]{MET-01}{El plan de comunicaciones cubre cinco grupos y no articula los trece actores de S2 ni define destinatario, frecuencia o canal. (rev. final, pp. 21, 23)}{S6 cruza los trece actores de S2 con destinatario, emisor, información, frecuencia, canal y escalamiento. T-9 resume las comunicaciones.}{S6, interesados, y T-9. Estado: Incorporada}
\observacion[sub=6]{MET-02}{Falta identificar quién prioriza y acepta incrementos por Curimón y quién lo reemplaza. (rev. final, pp. 21, 23)}{Curimón designa titular y suplente por acta antes del primer sprint. S6 y T-9 fijan atribuciones, plazos propuestos y escalamiento.}{S6 y T-9. Estado: Incorporada en diseño (designación contractual pendiente)}
\observacion[sub=6]{MET-03}{T-9 y T-10 se limitan a remisiones, con poca información utilizable por sí mismos. (rev. final, pp. 21--22)}{T-9 detalla control, gobierno y adhesión. T-10 incluye secuencia de sprints, capacidad, pruebas, criterios de salida y promoción de artefactos.}{T-9 y T-10. Estado: Incorporada}
\observacion[sub=6]{MET-04}{Los ocho sprints de A09 y cinco de A20 no tienen un plan de entregas por contexto o requerimiento. Sólo se detalla G13 en el sprint 7. (rev. final, pp. 22--23, 26--27, 48)}{Se planifican ocho sprints A09 y cinco A20 con requisitos, incrementos, dependencias, casos y perfil revisor. Se alinean G13 y los candidatos H4/H9 con la malla de S7. Se calcula capacidad neta con T-15 al planificar cada sprint.}{S6 y T-10. Estado: Parcial (nivelación por sprint con S7)}
\observacion[sub=6]{MET-05}{La deuda técnica se descuenta en sprints, pero carece de paquete y horas durante la operación Kanban. (rev. final, pp. 22--23)}{El carril Kanban 13.4 usa la capacidad variable de software durante M21 a M56. Deuda, mejoras y reservas en la misma bolsa se concilian con T-15 sin duplicar horas de cierre.}{S6 y T-10. Estado: Parcial (validar capacidad neta con S11 y cobertura del cierre en H1)}
\observacion[sub=6]{MET-06}{La viñeta de ramas temáticas de GitFlow cita normas ISO que no sustentan esa práctica. (rev. final, p. 23)}{El control de ramas se explica con documentación de ramas protegidas de GitLab.}{S6, ceremonias técnicas. Estado: Incorporada}
\observacion[sub=6]{MET-07}{Las tablas 6.1, 6.2, 6.3 y 6.6 contienen prosa. La Tabla 6.2 utiliza una descripción arquitectónica que no coincide con S4. (rev. final, p. 23)}{Se conservan datos comparables en las celdas y las reglas de cambio y aceptación en el texto. Se usan los seis contextos de S4.1.}{S6, comités, cambios, calidad y G13. Estado: Incorporada}
\observacion[sub=6]{MET-08}{[Observación complementaria Rev. 2] El S6 no menciona la adhesión de los 148 transportistas ni la trata como condición del método, aunque S7 le asigna paquetes EDT 9.1--9.3. (rev. final, p. 23)}{S6 y T-9 asignan al Líder Funcional el frente EDT 9.1 a 9.3, con control por RUT, cadencia, umbral de refuerzo y escalamiento.}{S6, interesados, T-9 y T-14. Estado: Incorporada}

\grupoObservaciones{EDT, cronograma e implantación (S7, T-14, T-15 y T-18)}
\observacion[sub=7]{PLA-01}{El primer nivel de la EDT se justifica genéricamente, divide el software por etapa y carece de agrupación intermedia. (rev. final, pp. 24, 27, 48)}{Los contextos y productos conservan código EDT; la lectura por capacidades cruza ambas etapas y los sprints detallan su entrega y dependencia.}{S7 capacidades y EDT; T-14; T-15 nodos A09 y A20 desagregados. Estado: Incorporada documentalmente}
\observacion[sub=7]{PLA-02}{Las láminas EDT repetían el diccionario en vez de mostrar la jerarquía. (rev. final, pp. 24, 27)}{Se regeneraron árbol y vistas jerárquicas desde el mismo catálogo de 54 paquetes; el diccionario queda en T-14.}{S7 anexos EDT y T-14. Estado: Incorporada}
\observacion[sub=7]{PLA-03}{Los cinco paquetes de innovación carecían de meses concretos en T-14 y de actividad identificable en la red. (rev. final, pp. 24, 27, 37)}{Cinco innovaciones con ventanas y nodos identificables, atribución de demanda reservada dentro de la capacidad mensual, sin duplicar nómina.}{T-14 EDT12; T-15 asignación por tarea, rol y fase; anexos vigentes. Estado: Incorporada documentalmente}
\observacion[sub=7]{PLA-04}{El esfuerzo se calcula desde dotaciones mensuales y no por tarea y paquete. (rev. final, pp. 25, 27, 48)}{Puente mensual por tarea, paquete, perfil y fase concilia exactamente con 99.626,67 horas internas de implementación y 84.606 de operación; se separan demanda directa, guardias y saldo de capacidad reservada.}{T-15 curvas y asignación; registro mensual de control. Estado: Incorporada como presupuesto de ejecución sujeto a revisión en H1}
\observacion[sub=7]{PLA-05}{A09 y A20 agrupan toda la construcción de cada etapa sin secuenciar los contextos dependientes. (rev. final, pp. 25, 27, 48)}{A09 se divide en ocho sprints y A20 en cinco con dependencias por capacidad; se recalculan ES, EF, LS, LF, HT y HL sin atribuir una ruta única a todos los nodos con HT cero.}{T-15 malla detallada; S7 red y Gantt de control. Estado: Incorporada}
\observacion[sub=7]{PLA-06}{La migración histórica y el estudio de costo por ruta no tienen precedencias ni holguras en la red. (rev. final, p. 25)}{El estudio de costo, la migración histórica, la preparación del cierre y el retiro del legado tienen nodos, límites, predecesoras y recursos propios.}{T-15 COST-M6, A08-HIST1, A24-HIST2 y CE2; T-14 EDT2.5/7.2. Estado: Incorporada}
\observacion[sub=7]{PLA-07}{La Gantt por paquete no muestra dependencias ni la ruta crítica de manera conciliable con T-15. (rev. final, pp. 25, 27, 48)}{Carta por paquete incorporada a T-14 y malla completa con dependencias y holguras en T-15; la vista CPM muestra la cadena de control de H12.}{T-14 carta Gantt; T-15 malla; S7 diagrama CPM. Estado: Incorporada documentalmente}
\observacion[sub=7]{PLA-08}{La ruta crítica encadena levantamiento, arquitectura, infraestructura, construcción y certificación con escaso paralelismo. (rev. final, p. 25)}{La red libera diseño preparatorio y ambientes de desarrollo para paralelizar trabajo sin esperar actas no habilitantes; conserva revisiones de aceptación y los doce hitos en su mes.}{T-15 relaciones FS/SS y fechas de liberación. Estado: Incorporada como línea base de planificación}
\observacion[sub=7]{PLA-09}{Las láminas Gantt de pocas filas y las 42 páginas gráficas fragmentaban la lectura. (rev. final, pp. 25, 27)}{EDT separada de la Gantt; doce vistas temáticas y una malla tabular sustituyen las conclusiones de la red anterior.}{S7 anexos activos y T-14. Estado: Incorporada}
\observacion[sub=7]{PLA-10}{La primera marcha blanca coincide con la temporada de fruta sin declarar la disponibilidad de usuarios de torre y terminales como dependencia. (rev. final, p. 26)}{Se reserva presencia en los 35 relevos semanales con cupos, reemplazo, simulación en PREPROD y certificación por perfil (S7/T-18); R-32 en S8/T-16 cubre formalmente la disponibilidad de usuarios en temporada de fruta.}{S7 partición de capacidad, T-18 capacitación y S8/T-16 R-32. Estado: Incorporada (agenda operativa a confirmar con terminales)}
\observacion[sub=7]{PLA-11}{A18 combina integración telemática en M13--M15 con montaje en M16--M18, pese a la restricción estacional. (rev. final, pp. 26--27)}{Se eliminó el nodo mixto A18: montaje físico en A11-PIL y M7/8/9-KIT; escalamiento lógico en A18-INT M13--M15.}{T-15 malla; T-14 EDT4.2/4.3; Gantt de plataforma. Estado: Incorporada}
\observacion[sub=7]{PLA-12}{El piloto A11 coincide con meses de cierres de Los Libertadores. La holgura se justifica con un cierre individual de hasta 12 días. (rev. final, p. 26)}{Los lotes tienen límites por ventana con agenda y stock previo (S7/T-18), y R-34 en S8/T-16 cubre cierres sucesivos de Los Libertadores sobre las ventanas fijas (días 112 a 180) eligiendo camiones que no cruzan frontera y reemplazando dentro del mismo lote.}{S7 calendario, T-18 y S8/T-16 R-34. Estado: Incorporada}
\observacion[sub=7]{PLA-13}{El plan de montaje no incorpora las 210 balizas de semirremolque, las 21 de reposición ni los lectores de portería declarados para la innovación 3. (rev. final, pp. 27, 48)}{El montaje distingue 210 balizas iniciales, 21 de reserva, cinco lectores y 19 equipos a bordo de reposición; banco, puesto, analista y servicio externo se imputan una vez.}{S3 balizas; T-14 EDT12.3; T-15 servicio físico y asignación; T-18. Estado: Incorporada en planificación}
\observacion[sub=7]{PLA-14}{[Observación complementaria Rev. 2] El Formulario T-14 no contiene la Carta Gantt (la remite a la sección 7.3.2 de S7) y la Etapa 2 carece de actividad de estabilización post-producción equivalente a A21 de Etapa 1. (rev. final, p. 27)}{T-14 incorpora directamente las doce láminas de Gantt. El cierre y estabilización E2 M21--M24 están separados de H12 y conciliados con la capacidad de soporte.}{T-14 carta; T-15 CE2 y cierre; T-18 M21--M24. Estado: Incorporada documentalmente}

\grupoObservaciones{Riesgos (S8 y T-16)}
\observacion[sub=8]{RIE-01}{La mitad del registro del cliente agrupa riesgos internos de audIT. R-06 trata como incertidumbre un retiro anunciado y ya resuelto por diseño. (rev. final, pp. 28, 31, 48)}{Cada riesgo declara su origen (mandante, terceros o audIT) y R-06 salió del registro activo. Son 35 riesgos: 17 de terceros, 12 de audIT y 6 del mandante.}{S8 §8.1 y T-16. Estado: Incorporada}
\observacion[sub=8]{RIE-02}{La convivencia y el retiro del sistema de 2013 no aparecen en el registro de riesgos. (rev. final, pp. 29, 31, 48)}{Queda como R-29, con retiro por función en M16, M21 y M24, conciliación diaria, reversión probada y acta del mandante. Su contingencia está en el S8.}{S8 §8.2 y T-16 R-29. Estado: Incorporada}
\observacion[sub=8]{RIE-03}{Faltan en el registro los riesgos de adopción de las innovaciones 1 y 2 y su riesgo residual. (rev. final, pp. 29, 31)}{R-30 (9 a 6) y R-31 (6 a 4), conciliados con las fichas 1 y 2 del T-19.}{S8, T-16 y T-19. Estado: Incorporada}
\observacion[sub=8]{RIE-04}{El árbol de fallas de asignación a conductor no apto carece de cuantificación. (rev. final, pp. 29, 31)}{Se fija el protocolo de conteo sobre unas 16.000 asignaciones del paralelo M13--M14, con intervalo binomial del 95 \% y la regla del tres cuando no hay eventos.}{S8 §8.2, árbol de fallas. Estado: Incorporada}
\observacion[sub=8]{RIE-05}{El análisis costo-beneficio de respuestas no presenta una razón de retorno. (rev. final, p. 30)}{Se define la razón de retorno de cada respuesta material como pérdida esperada evitada sobre costo, valorizada con las multas del Art. 80. Los montos van en el Sobre N.º 3.}{S8 §8.3. Estado: Incorporada}
\observacion[sub=8]{RIE-06}{R-17 se acepta con seguimiento, aunque la simulación reporta 21,6 \% de atraso de H1 y 16,1 \% de H2. (rev. final, pp. 30--31, 48)}{Sobre la malla conciliada, H4, H7, H11 y H12 quedan sin margen. R-17 pasa a exposición alta, con umbrales, escalamiento, recuperación con refuerzo y contingencia de entrega con alcance comprometido y acta. La malla se ajusta antes del H1 y la simulación se vuelve a correr. No se publican probabilidades de la red anterior.}{S8 §§8.2--8.3 y T-16 R-17. Estado: Incorporada (nueva simulación sujeta al ajuste de la malla antes de H1)}
\observacion[sub=8]{RIE-07}{Las tablas 8.1--8.3 se repiten entre S8 y T-16 y las remisiones son ambiguas. (rev. final, p. 31)}{Las tablas del T-16 se numeran T-16.n y el S8 remite a ellas por su número propio.}{S8 y T-16. Estado: Incorporada}
\observacion[sub=8]{RIE-08}{[Observación complementaria Rev. 2] R-22 trata suplantación digital pero no la desconexión física de antena/CAN/lector por conductores externos; la mitigación de R-14 es débil frente a 258 conductores externos al no apoyarse en el incentivo de la innovación 4 ni en el anexo de adhesión. (rev. final, p. 31)}{Nuevo R-35, desconexión física de antena, CAN o alimentación por un conductor externo, con exposición 12 a 8 y contingencia (equipo sin reporte en tramo con cobertura, sellos revisados en terminal, bloqueo de siguiente asignación y cláusula del anexo). R-14 se apoya en el anexo de adhesión y en el incentivo de comodato de la innovación 4.}{S8 §8.2 y T-16 R-14, R-22 y R-35. Estado: Incorporada}

\grupoObservaciones{Calidad y aceptación (S9, T-13 y T-17)}
\observacion[sub=9]{CAL-01}{El calendario de ensayos y las puertas de calidad se muestran sólo en tablas de detalle. (rev. final, p. 32)}{S9 incorpora la vista temporal vectorial (Figura 9.2) alineada con S7/T-15 y los hitos E-25 (actualizada en versión final con A20 en M13--M18 e hitos H3, H6, H8 y H11) y su explicación operativa.}{S9 §9.3.2, Figura 9.2 (PDF final verificado). Estado: Incorporada}
\observacion[sub=9]{CAL-02}{CP-HW-08 alude a polvo minero y CP-HW-06 duplica el ensayo de alarma SOS de CP-PERF-06. (rev. final, pp. 32, 34)}{CP-HW-08 fue saneado a estanqueidad IP67 en transporte terrestre (tanto en \texttt{catalogo-T-17.tex} como en \texttt{T-17.md}); CP-HW-06 verifica entrega SOS por evento con cobertura real y CP-PERF-06 mide P99 bajo degradación reproducible ($N \ge 100$).}{Formulario T-17 (\texttt{CP-PERF-06}, \texttt{CP-HW-06}, \texttt{CP-HW-08}). Estado: Incorporada}
\observacion[sub=9]{CAL-03}{Los 29 criterios de aceptación del capítulo 18 del caso no se vinculan a casos de prueba ni a actas. (rev. final, pp. 33--34, 48)}{La matriz \texttt{9-rtm-29} en T-17 traza los 29 criterios del Capítulo 18 a requisitos, casos de prueba, hitos y responsables; se cuadró además la trazabilidad de \texttt{CP-SYS-05} (\texttt{RF-002}) y \texttt{CP-PERF-03} (\texttt{RNF-014}, \texttt{RT-07.04}, \texttt{RT-07.07}).}{Formulario T-17 (Tabla \texttt{9-rtm-29} y catálogo). Estado: Incorporada}
\observacion[sub=9]{CAL-04}{La UAT de Etapa 1 coincide con el comienzo de la temporada de fruta, sin disponibilidad acordada de usuarios de terminal y despacho. (rev. final, p. 34)}{S9 §9.3.2 y T-13 §5 establecen ventanas UAT acotadas por turno fuera de horario punta, cupos nominales por rol (torre de despacho, prevención, mantenimiento, finanzas y terminales), suplentes designados y escalamiento en 48 h ante indisponibilidad para preservar H5; además, los simulacros semestrales de DR se fijan en junio y noviembre fuera de temporada alta.}{S9 §9.3.2 y Formulario T-13 §5. Estado: Incorporada}
\observacion[sub=9]{CAL-05}{El S9 ubica A12 en M10--M12 y el T-13 la expresa como M12/H5. (rev. final, p. 34)}{S9 y T-13 distinguen la ventana de ejecución de A12 (M10--M12) del hito de certificación H5 en M12, en total cuadratura con T-15 y E-25.}{S9 Tabla 9.5 y T-13 apartado 5. Estado: Incorporada}
\observacion[sub=9]{CAL-06}{El T-17 incluye un protocolo general y carece de protocolos por cada hito y del producto final. (rev. final, pp. 34, 48)}{El Formulario T-17 formaliza 13 fichas de protocolos individuales (H1 a H12 y Producto Final) alineadas a los Arts. 18.1--18.3 y garantía de correcto funcionamiento en M21 (Art. 37).}{Formulario T-17, protocolos de hitos H1--H12 y cierre. Estado: Incorporada}
\observacion[sub=9]{CAL-07}{[Observación complementaria Rev. 2] Las pruebas de integración con plataformas GPS de terceros (CP-INT-07 y CP-INT-08) se ejecutan sólo contra dobles simulados; ningún caso de staging o UAT define la muestra autorizada real con las tres plataformas. (rev. final, p. 34)}{Se incorporó en S9 §9.3.2, T-13 §5 y \texttt{CP-UAT-09} de T-17 la validación obligatoria en staging/UAT sobre muestra autorizada real de al menos una unidad activa por cada una de las tres plataformas GPS de terceros (SUP-02) previo al cierre de H5/H10, complementando los dobles simulados de CI (\texttt{CP-INT-07} a \texttt{CP-INT-09}).}{S9 §9.3.2, T-13 §5 y Formulario T-17 (\texttt{CP-UAT-09}). Estado: Incorporada}

\grupoObservaciones{Innovaciones y valorización (S13 y T-19)}
\observacion[sub=13]{INN-01}{La programación por oportunidad de patio tiene TRL 2 y metas relativas a una línea base sin medir. (rev. final, pp. 35--37, 45)}{La ficha fija la línea base en el piloto de M6, el método de medición, los umbrales (preparación inútil $-20\%$, primer intento $\ge 95\%$, sustitución $\le 30$ min) y la decisión del Comité al cierre de M6, con sustitución si no se acredita (contrastado en T-17 CP-UAT-02). El beneficio base queda en cero.}{S13 §13.2, ficha 2 del T-19 y T-17 CP-UAT-02. Estado: Incorporada}
\observacion[sub=13]{INN-02}{La innovación 3 presenta como novedad el kilometraje real de semirremolques, que apoya el criterio de aceptación 26. (rev. final, pp. 36--37)}{El kilometraje queda como alcance base del criterio 26. La innovación 3 se renombra «Acople verificado y registro térmico del semirremolque», igual que el paquete 12.3 del T-14.}{S13 §13.3 y ficha 3 del T-19. Estado: Incorporada}
\observacion[sub=13]{INN-03}{La innovación 5 presenta el catálogo de paradores del criterio 28 como parte del mérito innovador. (rev. final, pp. 35--37, 45)}{El catálogo es alcance base del criterio 28 y no se imputa a la innovación. La innovación 5 queda como «Anticipación circadiana y aviso no punitivo de descanso», igual que el paquete 12.5 del T-14.}{S13 §13.5 y ficha 5 del T-19. Estado: Incorporada}
\observacion[sub=13]{INN-04}{Las cinco fichas omiten el riesgo residual. La percepción de vigilancia en la innovación 5 se describe con probabilidad media a alta e impacto alto. (rev. final, pp. 36--37)}{Cada ficha declara su exposición inicial y residual enlazada al T-16: R-30 (9 a 6), R-31 (6 a 4), R-22 (8 a 4), R-11 (12 a 8) y R-14 (9 a 6). La innovación 5 añade disparadores de adherencia ($<70\%$ en M14) y NASA-TLX ($>35$).}{S13 y T-19, fichas 1 a 5. Estado: Incorporada}
\observacion[sub=13]{INN-05}{Las innovaciones 2, 3 y 5 carecen de base y magnitud para sus beneficios, aunque el informe remite la valorización al Sobre 3. (rev. final, p. 36)}{S13 explica el tratamiento de la planilla 3.2: beneficio cero en escenario base para innovaciones 2 y 3, valor pendiente de Etapa 1 para 1 y 5, y recupero bruto referencial de la 4 separado de la venta de audIT. Cada ficha fija el indicador físico a medir en piloto.}{S13 introducción y fichas, y Sobre N.º 3 (v3.2). Estado: Incorporada}
\end{tablaObservaciones}
